\chapter{Introduction}\label{chap:1}

Solving a Tangram together involves continually relating a small placement to a developing whole. One participant may try a triangle as a head while the other works on the body. A promising arrangement can guide the next move, expose a disagreement or be dismantled when the remaining pieces do not fit. The partners must judge the shape they are building while deciding how to respond to each other's actions. The view of the workspace is therefore part of the practical organisation of the task.

Shared visual information can help people follow task progress and establish what a partner means \cite{gutwin2002,gergle2013}. In virtual reality (VR), designers can change both the viewpoint and the way a participant reaches and manipulates objects. A first-person view places the user within the scene. An external view can expose the arrangement from outside it. For a planar puzzle, a perpendicular overhead view also provides a direct view of the outline and the relations between pieces. These properties make shared overhead viewing a useful setting for examining how visual access and collaborative action work together.

The project began with a supervisor-defined brief about collaborative behaviour in third-person VR. Its focus was how participants communicate intentions, coordinate actions and positions, negotiate object manipulation and maintain awareness of a shared workspace. Physical activity and conventional first-person VR provided comparison contexts. The implemented prototype develops this brief through a shared overhead view and a common Tangram task family. Both participants can manipulate pieces and decide how to divide or combine their work.

A shared view makes an arrangement available to inspect, but the pair still has to decide what to do with it. A participant can demonstrate an idea through a placement, continue a partner's partial solution or challenge it by trying another configuration. What remains on the table also matters: an earlier construction can become the starting point for later work. Following these sequences connects the display to the participants' use of it, including the moments when they keep a useful part, reopen it or accept a shape that does not match the target.

The research asks how this process develops in shared overhead VR and how comparison with other viewing conditions helps explain it. The overall research question is: How do pairs use shared overhead VR to judge and develop a Tangram construction together, and how do these processes relate to those observed in physical viewing and first-person VR?

RQ1: What shape-judgment problems arise during construction in the implemented viewing conditions, and how do participants check and respond to them?

RQ2: How do participants express local ideas through speech and manipulation, continue one another's work and reorganise who performs an action?

RQ3: How are partial constructions retained, extended or revised by their original builder and by the partner as the joint task develops?

Six pairs took part in physical direct viewing, physical overhead video, first-person VR and shared-overhead VR, in that order. The study uses researcher-prepared process records, group interviews, post-round questionnaires and available system records. The central analysis follows selected interaction sequences and reconstructs relevant VR arrangements from recorded piece poses. Action logs identify who controlled a piece and when; process records supply the surrounding proposal, response and task context. Together they make changes in both participation and geometry inspectable.

The findings bring three aspects of the task into one account. First-person VR included plausible constructions with unused pieces or incorrect proportions. G06's accepted candle provides a particularly clear example: incorporating all seven pieces resolved one visible problem while leaving a proportion mismatch. The pair's changing work on that construction also shows how initially separate local activity became shared reconstruction. Other cases show partners retaining or extending an arrangement that its original builder later changed, sometimes removing a correct partial solution.

The thesis contributes an empirical account of shared overhead VR within a broader set of collaborative viewing arrangements. Its contribution lies in connecting geometric judgment to the sequence of proposals, placements, continuations and revisions through which the pair produces a result. The fixed condition order, task patterns and prototype versions define the comparison and are considered in the method and discussion. Design implications concern access to the overall shape, usable manipulation, brief assistance and the preservation of earlier work during exploration.

Chapter~\ref{chap:2} establishes the concepts used to examine visible activity and collaborative problem solving. Chapters 3 and 4 describe the project's development, prototype and study method. Chapters 5 and 6 present the findings and discuss their implications for shared overhead VR. Chapter~\ref{chap:7} concludes the thesis.

